% GENERATED FILE. Edit canonical lessons, then run npm run generate:books.
% canonical-source-hash: fnv1a64:fdbe2fea3dc04b5b
% canonical-lessons: ES-C23-padre-madre, ES-C23-hermano-hermana, ES-C23-hermano-hache

\chapter{Parents and Siblings}
\label{ch:spanish-family}
\hlchaptermodality{68}

\begin{chapteropening}
\textbf{By the end of this chapter:} \emph{I can name my parents and my siblings in Spanish and explain why hermano is not built on Latin's word for brother.}

Give el padre, la madre, el hermano and la hermana, trace hermano back through germānus to germen, the seed, rather than to frāter, and separate its silent h from the once-spoken h of hijo and hacer.
\end{chapteropening}

\section[el padre, la madre]{el padre, la madre — among the oldest words there are}

\label{lesson:ES-C23-padre-madre}



\begin{quote}
\textquotedblleft{}Father\textquotedblright{} and \textquotedblleft{}mother\textquotedblright{} are some of the \textbf{oldest reconstructible words in human language} — older than Latin, older than Rome. Spanish kept them close to the Latin original, and they are secretly the \emph{same words} as English \textbf{father} and \textbf{mother}.
\end{quote}

\begin{cousinweb}
\begin{itemize}
\raggedright
  \item \textbf{el padre} (\textquotedblleft{}father\textquotedblright{}) $\leftarrow$ Latin \textbf{pater} $\leftarrow$ PIE \emph{\textbf{ph\textsubscript{2}t\'{\={e}}r}}.
  \item \textbf{la madre} (\textquotedblleft{}mother\textquotedblright{}) $\leftarrow$ Latin \textbf{māter} $\leftarrow$ PIE \emph{\textbf{méh\textsubscript{2}tēr}}.
\end{itemize}
\end{cousinweb}

\begin{culture}
\emph{Padre} and English \emph{father} both descend independently from the \textbf{same} PIE word — they're siblings, not parent and child. Latin (and Spanish after it) kept the original \emph{p} and \emph{t} largely intact; the Germanic branch that gave English took a separate road through \textbf{Grimm's Law}, which shifted the \textbf{inherited p $\to$ f} and \textbf{t $\to$ th}:

\noindent
\begin{tabularx}{\linewidth}{@{}>{\raggedright\arraybackslash}X>{\raggedright\arraybackslash}X>{\raggedright\arraybackslash}X@{}}
\toprule
\textbf{PIE ancestor} & \textbf{Latin $\to$ Spanish} & \textbf{Germanic $\to$ English} \\
\midrule
\emph{ph\textsubscript{2}t\'{\={e}}r} & \emph{pater} $\to$ \textbf{padre} & \textbf{f}a\textbf{th}er \\
\emph{méh\textsubscript{2}tēr} & \emph{māter} $\to$ \textbf{madre} & \textbf{m}o\textbf{th}er \\
\bottomrule
\end{tabularx}

So \emph{padre} and \emph{father} aren't lookalikes borrowed late, and neither one comes \emph{from} the other — they're the \textbf{same inherited PIE word}, one branch keeping the \emph{p}, the other shifting it to \emph{f}. The Latin root also gives Spanish and English their \textquotedblleft{}of a parent\textquotedblright{} adjectives: \textbf{paternal, paternidad, patrón} and \textbf{maternal, maternidad}. (A further Latin-only offshoot, \emph{māter} $\to$ \emph{māteria}, \textquotedblleft{}the stuff a tree-mother produces,\textquotedblright{} gives Spanish \emph{materia} and, separately, \emph{matriz} — \textquotedblleft{}womb\textquotedblright{} — via Latin \emph{mātrix}.)
\end{culture}

\subsection*{Your turn}
Say these aloud:

\begin{itemize}
\raggedright
  \item \textquotedblleft{}el padre, la madre\textquotedblright{}
  \item the sound-law — \textquotedblleft{}padre keeps p; father shifts to f — same word\textquotedblright{}
  \item \textquotedblleft{}padre $\to$ paternal; madre $\to$ maternal\textquotedblright{}
\end{itemize}

\begin{tcolorbox}[breakable,colback=teal!4,colframe=teal!35!black,title={Before you move on}]
Give \textquotedblleft{}father\textquotedblright{} and \textquotedblleft{}mother\textquotedblright{} with their articles. (\textbf{el padre, la madre}.) Does \emph{padre} turn into English \emph{father} directly? (\textbf{No} — both descend independently from the same PIE word; the Germanic branch shifted p $\to$ f and t $\to$ th via \textbf{Grimm's Law}, while Latin/Spanish kept the original sounds.) Name an English cousin of each. (\emph{padre} $\to$ paternal; \emph{madre} $\to$ maternal.)
\end{tcolorbox}
\section[el hermano, la hermana]{el hermano, la hermana — brother and sister, from a different word entirely}

\label{lesson:ES-C23-hermano-hermana}



\begin{quote}
You might expect Spanish \textquotedblleft{}brother\textquotedblright{} to come from Latin \textbf{frāter}, the way French's \emph{frère} does. It does not. Spanish's siblings took a completely different word — one about \textbf{seeds}.
\end{quote}

\begin{cousinweb}
Classical Latin could say \textbf{frāter germānus}, literally \textquotedblleft{}a brother \textbf{of the same stock/blood}\textquotedblright{} — distinguishing a \emph{full} brother from a half-brother, or from a \textquotedblleft{}brother\textquotedblright{} in some looser sense.

Then something ordinary and strange happened: in everyday Latin the \textbf{adjective} took over the whole job and the \textbf{noun quietly dropped out}. (English does this all the time — a \emph{mobile phone} becomes a \emph{mobile}, a \emph{cellular telephone} becomes a \emph{cell}.) What survived is the modifier:

\begin{itemize}
\raggedright
  \item \textbf{el hermano} (\textquotedblleft{}brother\textquotedblright{}) $\leftarrow$ Latin \textbf{germānus}.
  \item \textbf{la hermana} (\textquotedblleft{}sister\textquotedblright{}) $\leftarrow$ Latin \textbf{germāna}.
\end{itemize}

So Spanish's everyday word for brother is, grammatically, a leftover adjective.
\end{cousinweb}

\begin{culture}
\textbf{Germānus} traces back to \textbf{germen}, \textquotedblleft{}\textbf{sprout, bud, seed}\textquotedblright{} — the same root that gives English \textbf{germ}, \textbf{germinate}, and \textbf{germane} (originally \textquotedblleft{}closely related,\textquotedblright{} from the same seed, before it loosened into just \textquotedblleft{}relevant\textquotedblright{}).

So calling someone your \emph{hermano} is, at the etymological root, calling them \textquotedblleft{}\textbf{sprung from the same seed}.\textquotedblright{} (And no, this has nothing to do with \emph{Germany} — that name is a separate, unrelated puzzle.)
\end{culture}

\subsection*{Your turn}
Say these aloud:

\begin{itemize}
\raggedright
  \item \textquotedblleft{}el hermano, la hermana\textquotedblright{} — brother, sister
  \item the root chain — \textquotedblleft{}germen, seed\textquotedblright{} $\to$ \textquotedblleft{}germānus, of the same seed\textquotedblright{} $\to$ \textquotedblleft{}hermano\textquotedblright{}
  \item \textquotedblleft{}hermano\textquotedblright{} then English \textquotedblleft{}germinate, germane\textquotedblright{}
\end{itemize}

\begin{tcolorbox}[breakable,colback=teal!4,colframe=teal!35!black,title={Before you move on}]
Does Spanish \emph{hermano} come from Latin \emph{frāter}? (\textbf{No} — from \textbf{germānus}, \textquotedblleft{}of the same stock/seed,\textquotedblright{} an adjective that swallowed the noun \emph{frāter} it used to modify.) What does its root \emph{germen} mean, and what English words share it? (\textbf{Seed, sprout} — \emph{germ}, \emph{germinate}, \emph{germane}.) Next: the silent \emph{h} on the front, which is not the \emph{h} you think it is.
\end{tcolorbox}
\section[hermano from \emph{ermano}]{the \emph{h} in hermano — a letter that was never spoken}

\label{lesson:ES-C23-hermano-hache}



\begin{quote}
\emph{Germānus} starts with a \textbf{g}. \emph{Hermano} starts with an \textbf{h}. Something has to explain that — and the obvious explanation is the wrong one.
\end{quote}

\subsection*{What to know first}
\begin{itemize}
\raggedright
  \item el hermano, la hermana — the words and the \emph{germen} root.
\end{itemize}

\begin{culture}
You may already know that Spanish often turns a Latin initial \textbf{f} into a silent \textbf{h} — \emph{filius} $\to$ \textbf{hijo} (\textquotedblleft{}son\textquotedblright{}), \emph{facere} $\to$ \textbf{hacer} (\textquotedblleft{}to do\textquotedblright{}). \textbf{Hermano is not that}, and the difference is worth getting right.

\emph{Germānus} is stressed on its \textbf{third} syllable (\emph{ger-MĀ-nus}), which leaves the opening \emph{ger-} \textbf{unstressed}. Latin's \textbf{g} before a front vowel had already softened toward a \emph{y}-like sound in everyday speech — and in that unstressed position it simply \textbf{fell away entirely}, giving Old Spanish \emph{\textbf{ermano}}, with \textbf{no \emph{h} at all}. That is the form actually recorded centuries ago.

The \textbf{h} you see today was added \textbf{later, in spelling only}, to mark the word's old link to Latin \emph{g-}. Compare:

\noindent
\begin{tabularx}{\linewidth}{@{}>{\raggedright\arraybackslash}X>{\raggedright\arraybackslash}X>{\raggedright\arraybackslash}X@{}}
\toprule
\textbf{word} & \textbf{what happened to the sound} & \textbf{was the \emph{h} ever spoken?} \\
\midrule
\textbf{hijo}, \textbf{hacer} & Latin \emph{f-} $\to$ a real breathy \emph{h} & \textbf{yes}, for centuries \\
\textbf{hermano} & Latin \emph{g-} simply \textbf{dropped} & \textbf{no}, never \\
\bottomrule
\end{tabularx}

Same silent letter today; two completely different backstories, only one of which involves a sound at all.
\end{culture}

\subsection*{Your turn}
Say these aloud:

\begin{itemize}
\raggedright
  \item \textquotedblleft{}germānus\textquotedblright{} then \textquotedblleft{}ermano\textquotedblright{} then \textquotedblleft{}hermano\textquotedblright{} — the three stages
  \item the honest distinction — hijo/hacer's \emph{h} was once truly spoken; hermano's \emph{h} never was
\end{itemize}

\begin{tcolorbox}[breakable,colback=teal!4,colframe=teal!35!black,title={Before you move on}]
Is \emph{hermano}'s \emph{h} the same kind of change as \emph{hijo}'s or \emph{hacer}'s? (\textbf{No.}) What actually happened to \emph{germānus}'s opening? (Its \textbf{g} dropped away in an \textbf{unstressed} syllable, giving Old Spanish \emph{ermano}.) So where did the \emph{h} come from? (\textbf{Spelling} — added later to mark the Latin link, never pronounced.)
\end{tcolorbox}
